Another class of compounds formed by the iron-group elements consists of the so-called complex compounds.
Their distinguishing feature is that the transition-element atom does not enter the molecule as a simple ion.
Instead, it forms part of a composite, or complex, ion.
One example is the ion $\mathrm{MnO}_4^-$ in $\mathrm{KMnO}_4$.
Another is the ion $\mathrm{Fe(CN)}_6^{4-}$ in $\mathrm K_4Fe(CN)_6$.
In such complex ions the atoms

\SourcePage{381}{384}
lie closer together than in simple ionic compounds, and their $d$-electrons take part in valence bonding.
Accordingly, in complex compounds the iron-group elements behave like the elements of the palladium and platinum groups.
